\newpage % Rozdziały zaczynamy od nowej strony.
\section{Conclusions}

As a follow-up to the endpoint-prediction task, the following directions
could be pursued:

\begin{itemize}[label=\textendash]
  \item \textbf{Causal paths.}
  Greater emphasis could be placed on predictions along directed causal
  paths from drug exposure and patient context to clinical endpoints, together
  with a deeper analysis of what the model learns along those routes.

  \item \textbf{Real-world benchmarking.}
  Benchmarking the approach on MIMIC Clinical Database\footnote{%
    The MIMIC-III Clinical Database (v1.4) is a patient's dataset linking demographics, diagnoses, procedures, medications, and
    laboratory measurements that is used commonly in deep learning research.
    \url{https://physionet.org/content/mimiciii/1.4/}.%
  }
  could be a next step that would allow the pipeline to be validated on real clinical data.

  \item \textbf{GNN architectures.}
   Task~B was benchmarked with GraphSAGE, R-GCN, and TransformerConv; HGT was
    used in Task~A but not evaluated on endpoint prediction. A natural extension
    is to compare meta-relation attention (HGT) against TransformerConv on the
    same patient graph.

   \item \textbf{Statistical channel and HCR development.}
    Early experiments focused on pairwise dependence, while third-order evidence
    was implemented in a limited form, restricted to fully binary paths. Further
    work should devote more attention to triple-path evidence, and a dataset with
    richer numerical measurements at patient level would allow the Legendre
    expansion to be tested in larger scale.

  \item \textbf{Task~A--B loop.}
  Task~B currently uses a fixed graph; integrating it with Task~A link
  prediction would enable the closed-loop mechanism outlined in
  \cref{fig:loop}, in which graph updates and endpoint prediction inform one
  another iteratively.
\end{itemize}